\documentclass{article}
\usepackage[T1]{fontenc}
\usepackage{lmodern}
\usepackage{graphicx}
\usepackage{amsmath,amssymb}
\usepackage[a4paper,margin=2cm]{geometry}
\usepackage[absolute,overlay]{textpos}
\setlength{\TPHorizModule}{1mm}
\setlength{\TPVertModule}{1mm}
\usepackage[indonesian]{babel}
\usepackage{hyperref}
\setlength{\emergencystretch}{2em}
% unit: Halaman 25

\begin{document}

% === Halaman 25 (python/native) ===

\{ Enkripsi RC6-w/r/b

Masukan: Plainteks disimpan di dalam empat register berukuran

w bit: A, B, C dan D

r adalah jumlah putaran

S[0, ... , 2r + 3] adalah kunci internal, panjangnya w bit

Luaran: Cipherteks disimpan di dalam A, B, C, D

\}

Algoritma:

B B + S[0]

D D + S[1]

for i 1 to r do

t (B*(2B + 1)) <<< lg w

u (D*(2D + 1)) <<< lg w

A ((A t) <<< u) + S[2i]

C ((C u) <<< t) + S[2i + 1]

(A, B, C, D) = (B, C, D, A)

end

A A + S[2r + 2]

C C + S[2r + 3]

25

\end{document}
